%教程
\chapter[常用教程]{常用教程}

\section*{特别说明}
本节中所有上标“㊕”的网址仅在连接校园网时可成功访问相关服务。

\section[新生信息查询]{\uuline{新生信息查询}}
\label{freshman_query}
\begin{enumerate}
      \item 关注“山东第二医科大学学生之家”公众号
      \item 点击菜单栏中的“新生报到”→使用身份证号、手机号、验证码登录
      \item →根据其中的相关指引完成预报到流程
      \item →查看宿舍号、学号、班级、院系等相关信息，并根据个人需求选择是否订购军训套装、被褥套装\footnotemark 等
            \footnotetext{也可来校当场订购，详见系统及录取通知书说明。}
\end{enumerate}

\section[校园网]{校园网}
\subsection[无线连接]{无线连接}
\label{wifi_register}
\begin{enumerate}
      \item 激活：到校报到完成后，学校将按照个人身份证号在\uline{\href{http://210.44.80.65/}{http://210.44.80.65/}$^㊕$}为大家开通校园网账号，并告知初始密码\footnotemark
            \footnotetext{每年不同，详见班级群内具体通知。}
      \item \textbf{\uuline{充值}}：点击\uline{\href{https://slzfw.wfmc.edu.cn:8800/home/}{https://slzfw.wfmc.edu.cn:8800/home}$^㊕$}，并登录；或在校园网的登陆页面点击“自助服务”按钮\footnotemark，登陆后即可使用支付宝充值（微信支付暂时无法使用）
            \footnotetext{若点击按钮后网页显示“403 Error”，请将网址最前面的“http”改为“https”，按回车即可。}
      \item 校园网流量当前政策为每月免费30G，超出部分按0.5\textsf{￥}/G收费
\end{enumerate}

\subsection[有线连接]{有线连接}
\begin{enumerate}
      \item 按照\uline{\ref{wifi_register}} 的教程激活校园网
      \item 将网线插入宿舍内的“H3C”盒子的底部右侧的接口内并正确连接到电脑\footnotemark
            \footnotetext{若始终无法连接，应检查网线的内部线排列顺序，从左到右应为“白橙橙，白绿蓝，白蓝绿，白棕棕”。}
            \begin{enumerate}
                  \item 首次连接：
                        \begin{enumerate}
                              \item 打开“设置”→“网络和Internet”→“拨号”→“设置新连接”
                              \item →“连接到Internet”→点击“否，创建新连接(C)”→选择“宽带(PPPoE)(R)”
                              \item →输入自己的校园网帐号以及密码→勾选“记住密码”
                        \end{enumerate}
                  \item 再次连接：
                        \begin{enumerate}
                              \item 打开“设置”→“网络和Internet”→“拨号”
                              \item →选中之前设置的网络→连接即可
                        \end{enumerate}
            \end{enumerate}
      \item  注：上述均为Win10/Win11教程，Mac教程暂无。
\end{enumerate}

\section[校园手机卡]{校园手机卡}
\begin{enumerate}
      \item \textbf{开通与否全凭自愿，是否开通校园卡不影响校园网的使用。}随录取通知书一并寄出。如遇强制开通可告知带班学长或自行反馈
      \item 通常内含至少100分钟全国通话、80G校园流量（仅山东省内所有高校可用）\footnotemark
            \footnotetext{详情优惠政策可咨询营业厅，如追求更多流量建议学校各账号不绑定校园手机卡，每年根据新优惠政策调整手机卡（需及时注意注销旧手机卡以防欠费导致的信用记录问题）。}
      \item 开学报到当天可前往大服，免费领取礼品
\end{enumerate}

\section[空调使用教程]{空调使用教程}
\label{air_control}
\begin{enumerate}
      \item 微信关注“海享租”公众号，点击公众号菜单“在线租赁”，并注册、登录
      \item 点击“扫一扫”→扫描空调右下角二维码进行租赁\footnotemark
            \footnotetext{若提示租赁失败，请按照软件提示联系同宿舍的学长/学姐退租，也可咨询学长学姐或向宿管反馈。}
      \item →租赁完成后，点击“设备”→“空调图标”→“时长”，进行充值
      \item →点击“设备”→“空调图标”→“成员管理”，在此页面下将宿舍全部成员权限设置为均可管理空调开关即可
      \item \textbf{注意：}空调使用时长收费（0.55元/小时），具体收费及租赁政策详见“海享租”公众号
\end{enumerate}

\section[浴室预约与使用]{浴室预约与使用}
\subsection[浮烟山校区]{浮烟山校区}
\label{shower_software_f}
\begin{enumerate}
      \item 软件基础设置：
            \begin{enumerate}
                  \item 在手机应用市场下载“大白U帮”app
                  \item 按照实际住宿情况注册
                  \item 授予并开启“定位”与“蓝牙”权限
            \end{enumerate}
      \item 本楼层小浴室使用：
            \begin{enumerate}
                  \item 带好洗浴物品前往公共厕所旁边的浴室排队
                  \item 进入浴室，点击如右图所示的按钮（\mbox{\includegraphics[height=2.4ex]{bath.pdf}}）→选择“蓝牙设备”→“点击进行时”
                  \item →“洗澡”→“搜索洗澡”\footnotemark
                        \footnotetext{搜索不到设备请务必开启“蓝牙”功能，学校的设备无法扫码连接。}
                  \item →选择设备\footnotemark →“开始洗澡”
                        \footnotetext{距离厕所入口最近的是1号，远的是2号；不确定可以询问学长。}
                  \item 结束后点击“结束洗澡”按钮，并结算
            \end{enumerate}
      \item 一层公共大浴室预约：
            \begin{enumerate}
                  \item 在软件初始界面根据实际情况选择“X号楼1层”的浴室
                  \item →点击一个浴位，并点击“预约”按钮（若已满请选择“排队”）
                  \item →在8分钟内前往浴室，并点击“开始洗浴”
                  \item →结束后点击“结束洗澡”按钮，并结算
            \end{enumerate}
      \item 费用：以程序显示为准，详情收费标准略
      \item 申诉：如果在洗澡时突然停电导致无法结束洗澡而被扣费，请按照软件打开时弹出的公告，联系相关工作人员处理
      \item \textbf{注意：如果未点击“结束洗澡”按钮便直接离开可能会被多扣费}
      \item \textbf{严令禁止在浴室内大便！！！}
\end{enumerate}
\subsection[虞河校区]{虞河校区}
\label{shower_software_y}
\begin{enumerate}
      \item 支付宝搜索“住理生活”小程序，并按照提示开通账户（建议支付宝绑定银行卡）
      \item →搜索“潍坊医学院虞河校区”，并绑定账号到自己性别的浴室
      \item →扫描设备二维码即可使用淋浴
      \item 也可在小程序主页面点击“洗浴”→“切换设备”→手动选择设备进行洗浴
      \item \textbf{注意：浴室营业至晚10点，男浴室共14个浴位，请错峰洗澡}
\end{enumerate}

\section[洗衣机/洗鞋机使用教程]{洗衣机/洗鞋机使用教程}
\subsection[浮烟山校区]{浮烟山校区}
\label{washing_machine_f}
\begin{enumerate}
      \item 在微信小程序搜索“海乐生活”，并注册、开启相机与定位权限
      \item 预约方法（也可直接使用）：
            \begin{enumerate}
                  \item 在小程序内点击“附近营业点”→找到“潍坊医学院X号楼”
                  \item →选择相应的楼层→选中洗衣机并下单→等待上次洗衣结束
                  \item →前往洗衣机→输入验证码→放入衣物与洗衣粉/洗衣液并缴费
            \end{enumerate}
      \item 扫码直接使用（不可预约）：
            \begin{enumerate}
                  \item 前往洗衣机→在小程序内点击“扫码使用”按钮
                  \item →扫描洗衣机上的二维码→选择并下单
                  \item →放入衣物与洗衣粉/洗衣液，输入验证码后缴费即可
            \end{enumerate}
      \item 收费标准详见软件说明
      \item 洗衣机错误处理办法（若无相关经验切忌自行操作）：
            \begin{enumerate}
                  \item 拨打洗衣机旁边的报修电话；
                  \item E1：洗衣机断电后开门，打开洗衣机右下角小门，旋开阀门，使用镊子等工具伸入并清除其中堵塞管道的杂物，恢复原样即可；
                  \item E4：旋开洗衣机后方的水管阀门即可。
            \end{enumerate}
\end{enumerate}
\subsection[虞河校区]{虞河校区}
\label{washing_machine_y}
\begin{enumerate}
      \item 微信小程序搜索“智慧笑联”→按照提示注册并授予定位权限
      \item →绑定“潍坊医学院-虞河校区”→点击“洗衣机”
      \item →根据楼层选择洗衣机，按提示操作并付款即可
      \item 收费标准详见软件说明
\end{enumerate}
\subsection[禁止事项]{禁止事项}
\begin{enumerate}
      \item 禁止向第二格内倾倒洗衣粉、洗衣液，\textbf{第二格是放柔顺剂的}！
      \item \textbf{禁止使用洗衣机洗鞋}，请用旁边的洗鞋机！
      \item 禁止在洗衣机上堆放杂物
      \item \textbf{\uuline{禁止将袜子、内衣内裤等贴身衣物机洗！}}
\end{enumerate}

\section[烘干机使用教程]{烘干机使用教程}
\label{dry_machine}
\begin{enumerate}
      \item 注册等步骤详见\uline{\ref{washing_machine_f}}
      \item 按照预约的方法，选择宿舍楼后，选择“烘干机”即可，其他步骤与洗衣相似
      \item 推荐烘干配置
            \begin{enumerate}
                  \item 高温60分钟：大部分轻薄的衣物（例如T恤、卫衣、浴巾等）
                  \item 高温120分钟：薄被（如夏凉被）
            \end{enumerate}
      \item \textbf{注意：}\textbf{使用前后务必控干水箱并清理滤网。}棉被、羽绒服等禁止使用烘干机烘干以免损坏及不必要的危险情况发生。
      \item 收费标准详见软件提示
\end{enumerate}

\section[吹风机使用教程]{吹风机使用教程}
\label{hair_drier}
\begin{enumerate}
      \item 每层公共浴室旁边有两个公用吹风机，需扫码\footnotemark 租赁使用
            \footnotetext{部分吹风机屏幕二维码可能有缺损，不易扫描成功，多次尝试即可。（推荐使用浏览器扫码并复制到微信内收藏该网址，下次直接在微信内点击即可使用。）}
      \item →待手机发出“滴”-“滴”的声音后租赁成功
      \item →将手机扬声器对准吹风机租赁器方可正常使用
      \item 收费标准：详见软件提示，1分钱起步
\end{enumerate}

\section[图书馆座位预约教程]{图书馆座位预约教程}
\label{library_book}
\begin{enumerate}
      \item 微信小程序搜索“青栀校园”→微信注册登录并绑定学号→允许小程序通知
      \item →在小程序内点击“座位预约”或扫描图书馆座位上的二维码即可
      \item \textbf{座位暂离的注意事项：}
            \begin{enumerate}
                  \item 如因各种原因需要长时间离开的，请在小程序上选择“暂时离开”，否则按违规处理
                  \item 离馆时也需要在小程序内确认
                  \item 如发现已预约的座位被他人占据，请扫描桌面上的二维码并在小程序内举报，工作人员将尽快处理
            \end{enumerate}
      \item \textbf{违规说明：}
            \begin{enumerate}
                  \item 已预约而未按时到位的记一次违规
                  \item 未选择暂离而离开座位被举报的记一次违规
                  \item 停止使用后未选择退馆的记一次违规
                  \item \textbf{三次违规后将取消座位预约资格三天}
            \end{enumerate}
      \item 其他禁止事项
            \begin{enumerate}
                  \item 禁止占用其他人已经预约的座位
                  \item 禁止在图书馆内喧哗、吸烟
                  \item 禁止在图书馆的非背诵区域内背诵、朗诵、频繁交流
                  \item 禁止在图书馆内谈恋爱、亲嘴
            \end{enumerate}
\end{enumerate}

\section[设施报修方式枚举]{设施报修方式枚举}
\label{repair_report}
\begin{enumerate}
      \item 宿舍维修：
            \begin{enumerate}
                  \item 加入各宿舍楼的QQ报修群，在群内反映具体故障
                  \item 前往一层宿管处填表报修
                  \item 在宿舍一层宿管旁边的公告栏处查看相关负责人的电话，直接拨打即可
                  \item 拨打学生公寓管理中心电话反馈
                  \item 询问带班学长、学姐
            \end{enumerate}
      \item 教室维修：
            \begin{enumerate}
                  \item 拨打后勤管理处的电话报修
                  \item 拨打教室管理中心的电话报修（仅限多媒体及饮水机）
                  \item 拨打物业电话报修
                  \item 在“诉求留言”微信小程序内反馈
            \end{enumerate}
\end{enumerate}

\section[多媒体教室申请流程]{多媒体教室申请流程}
\begin{enumerate}
      \item 打印《多媒体教室使用审批表》并填写
      \item 前往本年级学工办签字、盖章
      \item 前往教室E区2层（见\uline{\ref{map_t}}）靠近A区处的“教室管理中心”签字盖章
      \item 前往预约的教室，拨打讲台上教室管理员的电话提前沟通说明
\end{enumerate}

\section[乐道济世书院教室申请流程]{乐道济世书院教室申请流程}
\begin{enumerate}
      \item 打印《入驻乐道济世书院申请表》并填写
      \item 前往本年级学工办签字、盖章
      \item 前往“书院管理办公室”（见\uline{\ref{map_a}}）签字盖章
\end{enumerate}

\section[学生会校级格式]{学生会校级格式}
\begin{enumerate}
      \item 用途：校级格式是一种用于校内各类\textbf{正式文稿}，如活动通知、综测条例、学生会文件等的标准
      \item 具体要求：
            \begin{enumerate}
                  \item 纸张大小：A4纸
                  \item 装订：页面左侧上下各 $\frac{1}{4}$ 处，距左边界0.3-0.5㎝处，钉与纸张左边界平行
                  \item 页边距：上下左右均为2.5㎝
                  \item 页码：两页及以上的材料，在页面底端居中插入
                  \item 行间距：固定值，26磅
                  \item 标点：均为中文标点，除特殊情况\footnotemark 外不应使用英文标点
                        \footnotetext{特殊情况举例：时间（22:05），英文活动（So, let us practice English!）等。}
                  \item 文章标题格式：方正小标宋简体，小二号，居中
                  \item 标题与正文之间空一行
                  \item 正文格式：仿宋-GB2312，三号，两端对齐，首行缩进2字符
                  \item 正文一级标题格式：“一、”；黑体，三号，居左，首行缩进2字符
                  \item 正文二级标题格式：“（一）、”；楷体-GB2312，三号，居左，首行缩进2字符
                  \item 正文三级标题格式：“1、”；其他要求与正文相同
                  \item 正文四级标题格式：“（1）”；其他要求与正文相同
                  \item 正文与落款之间空2行
                  \item 落款机构（个人姓名）格式：仿宋-GB2312，三号，右对齐
                  \item 落款时间格式：仿宋-GB2312，三号，右对齐，时间格式范例：2023年09月01日
            \end{enumerate}
\end{enumerate}

\section[小组汇报PPT制作指南（初级）]{小组汇报PPT制作指南（初级）}
\begin{enumerate}
      \item 整体要求
            \begin{enumerate}
                  \item 选择软件：请使用大多数人使用的微软\ Office办公软件，或者金山WPS最新版，并在电脑上进行编辑，手机当且仅当用于ppt的简单查看。如无特殊情况不要使用LibreOffice、OpenOffice、腾讯文档等办公套件以免不兼容。
                  \item \textbf{明确比例}：PPT（也称Slide）有4:3与16:9两种主流比例\footnotemark
                        \footnotetext{注：4:3比例像正方形，16:9是明显的长方形；下述的各类参数（如：24/32号字体）均以4:3、16:9的顺序进行。}
                  \item 首页规范：首页应含有所有的必要信息（详情见下文）
                  \item 目录规范：PPT应当在第二页含有一个简洁明了的目录
                  \item 内容规范：\textbf{PPT内仅应含有所讲内容的关键部分}而非一昧照抄原文
                  \item 字体与段落规范：\textbf{PPT字体不应过小，间距不应过密}（详情建议见下文）；禁止使用文字阴影；\textbf{特殊字体必须内嵌于PPT}（例如艺术字、书法体等），禁止现场在演示的电脑上下载补全需要的字体（操作步骤可搜索“在PPT内嵌入字体”）
                  \item 图片规范：每张图片不应大于10M\footnotemark，如发现图片模糊，需手动检查PPT设置，将PPT更改为“禁止自动压缩图片”，操作步骤请自行搜索
                        \footnotetext{如果图片过大，可以使用\uline{\href{https://meta.appinn.net/t/topic/44899}{Rimage\_GUI}}适当缩小图片体积（软件开源，如被报毒请自行分辨）。}
                  \item 动画规范：\textbf{PPT不应有过多动画以及元素堆叠}（例如，绝对禁止PPT中的一页内含15张大图片，依靠动画一张张切换，以免软件突然崩溃）
                  \item 配色规范：请\textbf{使用经典的“文字-背景”配色}（\sout{虽然确实难看}），例如“黑-白”，“红-白”，\linebreak[3]“白-黑”，“蓝-白”等，切勿使用“橙-白”、“红-蓝”等投影后效果一塌糊涂的配色\footnotemark
                        \footnotetext{若班内有色弱的同学，可适当换用该同学敏感、同时其他同学也能轻松接受的配色。}
                  \item 文件命名规范：要求PPT文件名\textbf{简单易懂，包含所有必须信息}
                  \item 文件保存规范\footnotemark：\textbf{必须同时以“.pptx”后缀与“.ppt”后缀各保存一份}以免部分电脑无法正常打开，\textbf{禁止保存为“.dps”、“.odp”等特殊格式}，详情见下
                        \footnotetext{\textbf{若无法看到文件名称后缀，请搜索“电脑设置显示文件扩展名”。}}
                  \item 学校图标使用规范：详情见\uline{\href{https://www.wfmc.edu.cn/4229/list.htm}{《山东第二医科大学VIS视觉识别系统手册》}}（由校宣传部印发）
            \end{enumerate}
      \item 详细要求：
            \begin{enumerate}
                  \item 首页：需注明小组成员姓名及学号，日期，课程名称，授课教师等必要信息并\textbf{仔细检查是否有误}，标题的字体不应小于80/75号，小组成员及其他信息字体不应小于25/25号
                  \item 目录页：简短凝练，总字数不应超过35个字；字体不应小于80/55号
                  \item 内容页：字体\footnotemark 不应小于55/40号，上下左右应各留出 $\frac{1}{10}$ 左右的间距；如内容过多应自行分页，\textbf{严禁为节省页数而缩小字号}；PPT显示的内容与口述补充的内容在6:4或7:3左右最佳，各类关键数据的引用（例如学术数据、课标外的公式定义等）应当按照\linebreak[3]《\textbf{GB/T}\ 7714—2015 信息与文献 参考文献著录规则》的相关标准\textbf{标明出处}
                        \footnotetext{下面的字体皆以微软雅黑为标准，楷体、仿宋等纤细字体请自行增加字号。}
                  \item 文件命名\footnotemark：推荐使用“\textbf{20XX级临床X班X组关于XXX的汇报（终稿）.pptx}”此类命名，\textbf{严禁使用默认名称以防混淆}（例如“新建 Microsoft PowerPoint 演示文稿.pptx”）。此外，在ppt未定稿时，推荐使用一些默认的规范进行命名（例如“\textbf{关于XXX的汇报草稿-4.pptx}”），以方便小组成员确定PPT版本，而非使用默认的“新建 Microsoft PowerPoint 演示文稿(1)（2）(5).pptx”这种高血压命名。
                        \footnotetext{文件名称后缀的“.pptx”、“.ppt”等各种类型并非手动添加，而是文件自带，\textbf{严禁手动修改后缀名！}}
                  \item 文件保存\footnotemark：如需在“.ppt”“.pptx”两种格式间相互转换请使用PowerPoint或WPS等办公软件，“.dps”格式必须使用WPS才能转换为“.ppt”或“.pptx”格式。
                        \footnotetext{\textbf{转换方法}：在左上角的“文件”菜单选择“另存为”，在下拉框中选择“.pptx”并保存。}
            \end{enumerate}

\end{enumerate}

\section[钉钉请假流程]{钉钉请假\footnotemark 流程}
\label{leave_dingtalk}
\footnotetext{需在学校统一将大家拉入钉钉的“山东第二医科大学”企业后方可使用。\textbf{各学院要求不一，仅以临床医学院为例。}}
\begin{enumerate}
      \item 线下请假步骤（正常情况）：
            \begin{enumerate}
                  \item 前往学工办或班主任办公室，当面请假并获得假条
                  \item →根据老师要求扫描相关二维码→钉钉填表
                  \item →刷脸进出校门，并将请假条之一交给保卫处
                  \item →返校后，在钉钉的电子假条处，以评论的方式销假
            \end{enumerate}
      \item 线上请假步骤：
            \begin{enumerate}
                  \item 打开钉钉→点击左上角选择主企业为“山东第二医科大学”
                  \item →点击页面最下方菜单栏“工作台”→点击“OA审批”
                  \item →在“学生日常事务管理”类选择“浮烟山校区本科学生请假单”→填表
                  \item →电话联系班主任老师或学工办老师，说明请假事由并等待审批
                  \item →审批通过后刷脸进出校门
                  \item →返校后，在钉钉的电子假条处，以评论的方式销假
            \end{enumerate}
\end{enumerate}

\section[家长、校友进校参观指南]{家长、校友进校参观指南\footnotemark}
\footnotetext{注：本节仅适用浮烟山校区；虞河校区为开放校区，除宿舍、教学楼外，均可参观。}
\begin{enumerate}
      \item 学校在每年8月30日前后允许家长参观校园\footnotemark，具体政策以学校官方说明为准
            \footnotetext{允许社会车辆在参观时段在校园内划定的区域内停放，（通常）允许家长在规定时间内参观宿舍环境，具体政策每年不同，具体情况以学校官方说明为准。}
      \item 其他校友若有返校需求可凭本人毕业证、学生证等有效证件（电子版也可），或通过当年所在院系的老师联系保卫处，在校门口登记后即可入校参观
\end{enumerate}

\section[缴费教程]{缴费教程}
\label{fee_pay}
\begin{enumerate}
      \item 官方：微信公众号“山东第二医科大学财务处”或“\uline{\href{http://tyzfpt.wfmc.edu.cn/xysf/login.aspx}{山东第二医科大学校园统一缴费平台}}”
      \item 用途：学费缴纳、卡号绑定等
      \item 学费缴纳教程：
            \begin{enumerate}
                  \item 前往公众号菜单，点击右下角“缴费管理”→“支付平台”
                  \item →登录系统（帐号为学号，初始密码为姓首字母大写加身份证后六位）
                  \item →按照提示修改初始密码（请务必牢记）→进行缴费
            \end{enumerate}
      \item 银行卡绑定教程：
            \begin{enumerate}
                  \item 目的：学校仅在初次使用时收集一次卡号并存储数据，以便下次直接使用\footnotemark
                        \footnotetext{详情见学校官方说明。}
                  \item 打开“山东第二医科大学财务处”公众号，点击“财务中心”
                  \item →使用帐号密码登录并绑定微信号（帐号为学号，初始密码为000000）
                  \item →点击“卡号维护”→“管理”
                  \item →按照提示填写相关信息，确认信息无误后提交即可
            \end{enumerate}
\end{enumerate}

\section[学工系统（微信小程序）]{学工系统（微信小程序）}
\begin{enumerate}
      \item 用途：学工系统主要用于晚点名、返校信息填报等日常工作\footnotemark
            \footnotetext{原“请假审批”、“外出审批”工作已基本转移至“钉钉”，教程参见\uline{\ref{leave_dingtalk}}。}
      \item 使用方式：
            \begin{enumerate}
                  \item 打开“定位权限”并允许微信使用→在微信搜索“智慧学工”小程序
                  \item 根据学校下发的账号密码进行登录（推荐立即与微信绑定以免忘记密码）
            \end{enumerate}
\end{enumerate}

\section[教务系统]{\textbf{\uuline{教务系统}}}
\begin{enumerate}
      \item 官网：\uline{\href{https://jwgl.wfmc.edu.cn/}{https://jwgl.wfmc.edu.cn/}$^㊕$}
      \item 用途：\textbf{选课，缓考申请，成绩查询}，查看（导出）课程表，空闲教室查询
      \item 缓考申请教程：
            \begin{enumerate}
                  \item 进入教务系统
                  \item 点击左侧菜单“考试报名”→“我的申请”→“缓考申请”
                  \item →选择“学年学期”和“活动名称”后，直接点击“搜索”（不要填写科目名称）
                  \item →在弹出的菜单中选择缓考科目并填写申请\footnotemark
                        \footnotetext{若因病缓考体测，需要上传病历本等相关材料，并前往校医院开具证明，再前往学工办向教师当面说明情况。}
            \end{enumerate}
      \item \textbf{注意：}仅限校内访问，如需在家使用教务系统，参见\uline{\ref{cas_system}}条目
\end{enumerate}

\section[资源访问控制系统（校内VPN）]{\textbf{\uuline{资源访问控制系统（校内VPN）}}\footnotemark}
\footnotetext{在校内请直接访问相应系统，无需使用本系统中转。}
\label{cas_system}
\begin{enumerate}
      \item 官网：\uline{\href{https://webvpn.wfmc.edu.cn/}{https://webvpn.wfmc.edu.cn/}}
      \item 说明：（也称CAS认证系统）本系统用于在校外访问校内网络信息资源，如：教务系统（查成绩、课表）、知网、临床医学虚拟仿真实验中心\footnotemark 等
            \footnotetext{查阅文献推荐使用CARSI系统，速度快效果好，教程详见\uline{\ref{carsi_system}}。}
      \item 异地登录教务系统教程（其它系统同理）：
            \begin{enumerate}
                  \item 打开网站，点击“统一身份认证登录”，按照学校下发的专用账号、密码登录即可（推荐绑定微信）
                  \item →找到应用中心→“教务系统-非单点登录”\footnotemark →使用教务系统账号密码登录即可
                        \footnotetext{请注意，在校外时点击“教务系统”无法登录，只有“非单点”能校外登录！}
            \end{enumerate}
\end{enumerate}

\section[CARSI系统]{\textbf{\uuline{CARSI系统}}}
\label{carsi_system}
\begin{enumerate}
      \item 官网：\uline{\href{https://ds.carsi.edu.cn}{https://ds.carsi.edu.cn}}
      \item 说明：（与CAS认证系统不是同一个系统）用于快速访问学校订阅的各类数据库，如百度文库、知网、万方、维普等
      \item 使用教程：
            \begin{enumerate}
                  \item 进入官网→搜索“潍坊医学院”并勾选“记住我的选择”→点击后进入登陆界面\footnotemark
                        \footnotetext{请注意，不要收藏登录界面的网址！每次登录网址都不一样，只能从官网重新进入！}
                  \item →使用\textbf{校园网账密}登录系统，出现各类弹窗一律选择“Accept”即可\footnotemark
                        \footnotetext{仅推荐在自己的电脑上如此设置，如必须在网吧等公共场所的电脑上使用，请审慎阅读相关提示，并谨慎进行登录，如因账号泄露造成损失，一切责任自负。}
                  \item →登录完成后，点击任意资源链接即可进入相应网站并获取论文
            \end{enumerate}
\end{enumerate}

\section[校务行（微信小程序）]{校务行（微信小程序）}
\label{cert_prover}
\begin{enumerate}
      \item 官网：微信小程序
      \item 用途：查成绩，下载学籍证明、成绩证明的pdf版本
      \item 费用：以程序显示为准
      \item 教程：
            \begin{enumerate}
                  \item 搜索小程序“校务行”
                  \item →点击右上角“点击登录”（帐号为学号，密码为身份证后六位）\footnotemark
                        \footnotetext{新生可能无法在开学后即刻使用本程序，需等待学校将信息录入完毕。}
                  \item →按需选择“电子成绩单”或“电子证明”→按照小程序提示进行即可
            \end{enumerate}
\end{enumerate}

\section[常见证明申请]{常见证明申请}
\begin{enumerate}
      \item 说明：适用于学籍证明、成绩证明
      \item 线下打印
            \begin{enumerate}
                  \item 携带身份证、手机前往D区的打印机，详情位置见此\uline{\ref{common_locations}}
                  \item 按照机器的说明填写相关信息并打印
                  \item 无纸时请按照机器上的相关联系方式联系工作人员
            \end{enumerate}
      \item 线上下载
            \begin{enumerate}
                  \item 按照\uline{\ref{cert_prover}}的教程申请并下载，自行前往打印店打印即可（请注意及时删除打印完毕的文件以免泄漏隐私）
            \end{enumerate}
\end{enumerate}

\section[高铁学生票购买流程]{高铁学生票\footnotemark 购买流程}
\footnotetext{仍保留线下优惠资质核验、学生票购买渠道，若操作遇到问题可直接前往线下售票处，通过工作人员核验与购买；详细的注意事项等其他相关内容参见12306官网或其官方app。}
\begin{enumerate}
      \item 等待学校完成学籍注册并下发加盖注册章的学生证，确认学生证上有“火车票学生优惠卡”\footnotemark
            \footnotetext{若学生证丢失请以学籍证明替代并在人工窗口办理相关业务；学生证每学期以班级为单位统计一次并统一补办。}
      \item 首次使用请先在\uline{\href{https://www.chsi.com.cn/}{学信网}}查看学籍，确认信息已更新
      \item 打开12306 app→进入“我的”页面→点击“学生优惠资质核验”旁的“点击查看”按钮
      \item →填写学生资质信息→等待审核结果（3个工作日内反馈）
      \item 价格：高铁75折，普通火车5折（仅二等座可使用优惠，详见相关规定）
      \item \textbf{注意：}每学年\footnotemark 有4次优惠机会；每学年需重新核验一次优惠资质
            \footnotetext{每年10月1日至下一年的9月30日为一个学年。}
\end{enumerate}

\section[文体中心预约教程]{文体中心预约教程}
\label{sports_center_book}
\begin{enumerate}
      \item 关注“潍医文体中心”公众号→点击“场地预约”→“预约入口”
      \item →点击“我的”→“校内登录”→使用CAS认证系统的账号密码登录
      \item 进入个人中心页面，点击“人脸录入”，录入信息→
      \item 按需预约游泳馆、羽毛球馆等，并付费→
      \item 在预定时间段\footnotemark 前往场馆，向工作人员出示二维码即可
            \footnotetext{场馆开放时间见\uline{\ref{sports_center_operating_hours}}。}
      \item 费用：
            \begin{enumerate}
                  \item 健体中心：3元/2小时
                  \item 羽毛球馆\footnotemark：6元/时/片 3元/人/2小时
                        \footnotetext{羽毛球等多人运动项目预约方式为“单人预约，多人共享”，详情可咨询负责相关场馆的教师。}
                  \item 篮球馆：20元/1小时（半场）或40元/1小时（全场）
                  \item 游泳馆\footnotemark：8元/场/2小时
                        \footnotetext{注意：游泳馆仅限本人购票，代购无效，每人每场限购一张。}
                  \item 室内乒乓球场：2元/台/2小时
                  \item 室内网球场：6元/片/1小时
            \end{enumerate}
      \item 特殊说明：因羽毛球馆和篮球馆共用同一场地，故此二者互斥；场馆开放具体时间以及价格变动以公众号通知为准
\end{enumerate}
